\documentclass[../main.tex]{subfiles}
\begin{document}
%==========================================TIÊU ĐỀ TẬP========================================
\begin{table}[h]
    \begin{adjustwidth}{-1cm}{}
        \begin{tabular}{>{\centering\arraybackslash}p{6cm}|>{\raggedright\arraybackslash}p{14cm}}
            \multirow{5}{6cm}
            % Cột bên trái: ÉPISODE và số tập
            {\\[-40pt]
                \flaregothic\fontsize{20pt}{20pt}\selectfont \centering
                \color{eptype}ÉPISODE\color{black}\\
                \burbank\fontsize{72pt}{72pt}\selectfont
                \color{epnum}3\color{black}\\[6pt]
                \cafeteria\fontsize{20pt}{20pt}\selectfont\color{epnum}(0-4)\color{black}
                \\[18pt]
                \flaregothic\fontsize{14pt}{14pt}\selectfont
                \color{epattr}BIENVENUE, \uline{\color{Haachama} Haato} et \uline{\color{Ojou} Ayame} !\\
                \color{black}
            }
             &
            % Dòng thứ nhất: tên tập tiếng Nhật
            {\fotskip\fontsize{18pt}{18pt}\selectfont \ruby{前}{ぜん}\ruby{回}{かい}のラブライブ！【ホロライブ】
            }                                                                             \\[6pt]
            % Dòng thứ hai: tên tập tiếng Anh
             & {\cafeteria\fontsize{18pt}{18pt}\selectfont Last Time, on \hll !}                                                                                        \\[4pt]
            % Dòng thứ ba: tên tập tiếng Việt
             & {\vnmsans\fontsize{18pt}{18pt}\selectfont Kỳ trước ở \hll!}                                                                                              \\[8pt]
            % Dòng thứ tư: Nhân vật xuất hiện
             & {\caviar{\fontsize{14pt}{14pt}\selectfont Nhân vật: \emp{kuon1} \emp{roboco} \emp{fubuki} \emp{haato} \emp{matsuri} \emp{aqua} \emp{ayame} \emp{subaru}}}
            \\[8pt]
            % Dòng cuối cùng: Thời gian diễn ra của tập (thường sẽ là ngày phát hành)
             & {\continuum\fontsize{16pt}{16pt}\selectfont Ngày phát hành: 19/05/2019}                                                                                  \\[18pt]
        \end{tabular}
    \end{adjustwidth}
\end{table}
%=========================================TRUYỆN CHÍNH========================================
\color{black}

\txtbx[2]{
    {\color{vol} \uline{Tóm tắt:}} Một chuỗi những sự kiện ngẫu nhiên và hỗn loạn được kể lại dưới hình thức hồi tưởng ``Kỳ trước ở \hll!''. Từ căn bệnh cười quái ác của Subaru, cho tới thảm họa dị ứng phấn hoa khiến Roboco hắt xì ra dầu ngập lụt cả văn phòng. Khó khăn lại chồng chất khó khăn khi Aqua trượt ngã đập đầu đến mức mất trí nhớ, và rồi mọi thứ hoàn toàn vượt ngoài tầm kiểm soát, kết thúc bằng một vụ hỏa hoạn thiêu rụi toàn bộ công ty!
}

\begin{center}
    (\uline{Những đoạn phim dưới đây được phát lại dạng ``Tóm tắt tập trước'' của series \textit{Love Live!} - tức \textit{Zenkai no Love Live!}})
\end{center}

\subsection*{\phantomsection\label{ep-003.I}}
\addcontentsline{toc}{subsection}{{\sen{-I-}} {\caviar Căn bệnh "nhìn cái gì cũng cười" của Subaru}}
\stpt{-I-}

\stcap{Căn bệnh "nhìn cái gì cũng cười" của Subaru}
\begin{center}
    \textbf{Kỳ trước ở \hll !}
\end{center}

\pov{Oozora Subaru}

\gif{0-4a}{1}{35}

Tôi chạy vội về hướng chiếc camera đang quay một vòng quanh phòng, bật nhảy lên thật cao và hét lớn: ``\textbf{Trong kỳ trước của \hll!}''. Đây chính là cách nhái lại màn xuất hiện kinh điển của Kousaka Honoka - trưởng nhóm $\mu's$ - trong phân cảnh mở đầu tập một của bộ anime huyền thoại \textit{Love Live! School Idol Project}.

Bất giác, tôi tự nhớ lại hành trình trưởng thành đầy gian nan của chính mình. Những hình ảnh cũ hiện lên sống động: tôi đứng cùng với gia đình vào những dịp tựu trường đầu tiên, từ những ngày chập chững bước vào trường mẫu giáo, rồi lên tiểu học, trung học cơ sở, và giờ là trung học phổ thông.

Thế nhưng, giờ đây tôi đã chạm đến độ tuổi ẩm ương dở dở ương ương, cái tuổi mà người ta bảo là ``chỉ nhìn thấy cọng đũa lăn thôi cũng đã cười''. Đến mức bây giờ, bất cứ thứ gì trên đời này cũng có thể khiến tôi bật cười nắc nẻ! Ôi trời đất ơi, tại sao tôi lại mắc phải cái hội chứng dậy thì tai quái này cơ chứ!?

\gif{0-4b}{1}{111}

Thật vậy đấy, chẳng đùa đâu! Chỉ cần một chiếc đũa vô tình lăn nhẹ trên mặt bàn ăn thôi là tôi đã cười khoái chí, cười ngặt nghẽo đến mức ngã lăn lộn từ trên ghế bành xuống đất rồi!

Anh Kuon ngồi đối diện nhìn tôi một hồi lâu, rồi anh ấy khẽ lắc đầu, ngao ngán cảm thán: ``Cái gì vậy trời\dots? Tự dưng lại cười sặc sụa như bị lên cơn ngáo đá thế kia?''

Tôi vừa cười sặc sụa vừa cố gắng đáp lại lời anh: ``Có phải em muốn thế đâu! Nhưng mà-- ahahaha!'' rồi lại tiếp tục ôm bụng cười ngặt nghẽo trước bộ mặt nhăn nhó, bất lực toàn tập của anh ấy.

Nhưng tôi tự nhủ bản thân không thể cứ tiếp tục sống buông thả như thế này mãi được! Cách duy nhất và triệt để nhất để chữa dứt điểm căn bệnh quái ác này là tuyệt đối không được nghĩ đến chuyện dậy thì nữa!

Trong tâm trí, tôi mơ màng tưởng tượng về một ngày mình có thể cao lớn vượt trội về mặt thể xác, sở hữu một gương mặt trưởng thành, nghiêm nghị cùng một dáng vóc oai phong lẫm liệt. Một vầng hào quang rực rỡ sẽ bao quanh lấy tôi, và tôi sẽ tự tin dõng dạc tuyên bố với cả thế giới rằng mình đã hoàn toàn khỏi bệnh!

% \img{0-4c}

``Bằng cách thề nguyền với thần linh, mình đã có thể lớn nhanh như thổi! Thế nhưng\dots''

Ôi, thực tế lúc nào cũng phũ phàng và tàn nhẫn dội một gáo nước lạnh vào mặt tôi. Một cô nàng Quỷ đang nằm chống cằm lười biếng trên ghế sô-pha, khẽ ngẩng đầu lên nhìn tôi rồi buông một câu nhận xét lạnh lùng khiến tôi lập tức `hóa đá': ``Thực tại không giống như trong manga đâu nghen.''

Chán thật đấy, nhưng mà người ta cũng có quyền được mơ mộng chút chứ bộ!

Chưa dừng lại ở đó, anh Tổng ngồi gần đó cũng gật đầu đồng tình, hùa theo để nhấn mạnh sự vô lý trong suy nghĩ của tôi bằng một lời khuyên thực tế: ``Đúng đó. Ý tưởng của em nghe viển vông quá sức rồi. Kiềm chế lại bản thân có khi còn hợp lý và khả thi hơn đấy.''

Trái lại, cô nàng hậu bối có mái tóc vàng óng đang ngồi trên chiếc ghế xoay bên cạnh lại tỏ vẻ thích thú trước ý tưởng bất khả thi nhưng đầy thú vị này: ``Chị lại thấy nhìn em như thế này cũng hay ho đấy chứ.''

\model{24}{r}{8cm}{ayame}

\chrint

Nàng Quỷ vừa cất lời đáp lại Subaru ban nãy chính là Nakiri Ayame (\ruby{百}{な}\ruby{鬼}{きり}あやめ, \ruby{Na}{Bách}-\ruby{ki-ri}{Quỷ}\ A-gia-mê), một thành viên vô cùng nổi bật thuộc Thế hệ thứ Hai. Cô là một cô nàng Oni (quỷ sừng) đang giữ chức danh Hội trưởng hội học sinh tại Học viện Ma giới đầy quyền lực, và luôn giữ thói quen nói chuyện với tông giọng bề trên đầy kiêu hãnh đối với loài người. Dù mang danh là quỷ, thế nhưng cô lại cực kỳ thân thiện, tốt bụng và luôn cư xử bình đẳng với tất cả mọi người xung quanh. ``Konnakiri!'' chính là câu chào thương hiệu của cô nàng mỗi khi bắt đầu các buổi phát sóng trực tiếp của mình.

Nàng Oni này sở hữu một khiếu hài hước cực kỳ nhạy bén, cô rất dễ bật cười với mọi thứ trên đời cho dù đó chỉ là những trò đùa vô cùng giản đơn, đặc biệt là những câu chơi chữ tiếng Nhật nhạt nhẽo. Đây cũng chính là điểm yếu chí mạng của Ayame, bởi nhiều lúc cô cười không thể kiểm soát nổi đến mức phải đập bàn liên tục để kiềm chế cảm xúc phấn khích của mình. Cô nàng cũng là một kẻ rất thích trêu chọc người khác, vì vậy việc bị Kuon gọi là một cô gái ``loi nhoi và thích đùa giỡn'' hoàn toàn chẳng sai chút nào.

Bên cạnh đó, Ayame cũng là một người có khả năng phân tích cực kỳ sắc bén về những diễn biến trong xã hội. Cô luôn chủ động tìm hiểu thật kỹ đối tượng trước khi quyết định cách ứng xử phù hợp. Với tư cách là một học sinh đến từ thế giới bên kia đang nỗ lực tìm hiểu cách nhân gian sinh sống và làm việc, đây quả thực là một điểm cộng vô cùng lớn.

Tuy nhiên, cô nàng cũng có những lúc khá nóng tính nếu các trò đùa bị đẩy đi quá trớn. Ayame từng đấm thủng cả tường nhà vì bực tức, thậm chí là đánh nhau với cả bố mình. Mặc dù cô luôn cố gắng kiềm chế cơn giận dữ và duy trì trạng thái vui vẻ, đáng yêu trước mặt những người hâm mộ của mình, nhưng bản chất Oni mạnh mẽ đôi khi vẫn khiến cô gặp không ít chật vật trong việc làm chủ cảm xúc của bản thân.

Ayame sở hữu một mái tóc dài màu trắng ánh hồng với hai búi tóc buộc gọn gàng ở phía sau, hai chiếc sừng trắng nhỏ nhắn nhô lên đầy kiêu hãnh trên trán, đi kèm một cuộn dây đen buộc ở gốc sừng bên trái. Đôi mắt đỏ hồng rực rỡ luôn lấp lánh niềm vui, kết hợp cùng chiếc chuông đỏ và mặt nạ hình oni trang trí trên tóc. Trang phục thường ngày của cô là bộ kimono màu đen viền vàng ngắn ngang đùi với phần tay áo rộng được mặc trễ vai quyến rũ, khoác bên ngoài là một bộ kimono không tay màu đỏ rực rỡ, được cố định bằng một sợi dây buộc xoắn màu trắng-đỏ nổi bật phía sau lưng. Cô đi một đôi tất dài màu trắng quá đùi cùng dép okobo đỏ cao gót, chỉ để lộ ba ngón chân giữa. Và đặc biệt, phía sau lưng cô luôn gác sẵn hai thanh katana sắc bén, dường như là để sẵn sàng ``trừng trị'' bất cứ kẻ nào dám chọc giận mình\dots\ (hy vọng rằng cô nàng sẽ chẳng bao giờ phải dùng đến chúng.)

\model{24}{l}{6.5cm}{haato}

Cô nàng vừa lên tiếng tiếp lời Ayame chính là Akai Haato (\ruby{赤}{あか}\ruby{井}{い}はあと, \ruby{A-ca}{Thích}-\ruby{i}{Tĩnh}\ Ha-tô), một thành viên đầy cá tính thuộc Thế hệ thứ Nhất. Cô là một nữ sinh trung học sở hữu tính cách láu lỉnh, đôi lúc cáu kỉnh và được mệnh danh là một trong những thành viên có tính cách dị biệt, lập dị nhất trong cái văn phòng này. Với tư cách là một kouhai (hậu bối), cô luôn nỗ lực tìm mọi cách để thu hút sự chú ý từ mọi người nhằm kết bạn, cho dù đôi khi chính tính cách có phần kỳ quái ấy lại là rào cản lớn khiến cô gặp khó khăn.

Haato là một tsundere điển hình, cô cực kỳ thích trêu chọc những người hâm mộ của mình và thường gọi họ là những chú ``lợn'' hay ``heo'' kinh tởm bằng một chất giọng đầy châm biếm, mặc dù thâm tâm cô hoàn toàn không hề có ác ý gì. Tính cách ``khẩu thị tâm phi'' của cô còn được thể hiện rõ qua việc cô là một trong những thành viên `trong sáng' nhất của đại gia đình \hll, dù đôi lúc cô cũng có những pha xử lý vô cùng táo bạo trong các buổi phát sóng trực tiếp. Bù lại, Haato là một người rất điềm tĩnh và biết lắng nghe, luôn sẵn sàng làm chỗ dựa vững chắc cho mọi người xung quanh tìm thấy sự bình yên.

\dots Hình như cô nàng còn sở hữu một nhân cách thứ hai nữa, nhưng có vẻ như cô ấy đã kiểm soát nó tương đối tốt rồi. Nghe đồn trong các buổi phát sóng, nhân cách bí ẩn đó xem chừng cực kỳ điên loạn và đáng sợ\dots\footnote{Về nhân cách thứ hai của Haato, tức Haachama, người viết xin phép sẽ không đề cập sâu trong cuốn truyện này.}

Đặc biệt, Haato có một khuyết điểm cực kỳ chí mạng là kỹ năng nấu ăn siêu dở tệ. Trong nhiều lần phát sóng trực tiếp, cô đã tạo ra những món ăn được ví như ``thảm họa ẩm thực không thể tiêu thụ nổi'', khiến cho tất cả mọi người xung quanh phải đặt một dấu hỏi lớn về trình độ của cô. Thế nhưng, cô nàng vẫn luôn cực kỳ tự tin, thậm chí là có phần cao ngạo về tài năng nấu nướng của bản thân mình.

Haato có một mái tóc vàng óng dài thướt tha, được buộc thấp thành hai chỏm ở hai bên bằng một sợi ruy-băng đỏ dài cùng dây chun trang trí hình trái tim đỏ thắm. Đôi mắt xanh nhạt của cô luôn toát lên vẻ tinh nghịch. Tên của cô có nghĩa là ``màu đỏ'' và ``trái tim'', vì thế cũng dễ hiểu khi cô nàng luôn thích đính những biểu tượng này lên tóc và trang phục của mình. Cô thường mặc một chiếc áo sơ-mi trắng ngắn tay kết hợp cùng chiếc đầm màu xanh cao cổ nhưng lại ngắn ngang đùi đầy năng động, đi kèm dây đeo lưng màu đỏ, đùi trái đeo một chiếc đai đỏ nổi bật, diện cùng đôi tất đen dài quá gối và giày lười học sinh truyền thống.

\chrint

Tôi ngay lập tức phản ứng lại, cố chấp bảo vệ đến cùng ý tưởng của mình: ``Nhưng em bắt buộc phải làm vậy thì mới chữa khỏi bệnh được chứ! Em không thể cứ chịu trận chịu đựng như thế này mãi được đâu!'' Dù có bị cả hai người kia thẳng thừng bác bỏ ý kiến ban nãy, nhưng sự ủng hộ từ một tiền bối lập dị như chị Haato lại càng tiếp thêm cho tôi vô vàn quyết tâm.

Tôi bày ra một vẻ mặt nghiêm túc tột cùng, nhìn thẳng vào Kuon-senpai với ánh mắt rực cháy ý chí chiến đấu. Thế nhưng, anh ấy chỉ đáp lại tôi bằng một cái nhìn đầy nghi ngờ và ái ngại, thầm tự hỏi trong đầu: ``\textit{Cái cô nàng Vịt này rốt cuộc lại định bày ra trò điên rồ gì tiếp theo đây\dots?}''

Cứ chờ đấy, rồi em sẽ có cách cho anh xem, Kuon-senpai!

\subsection*{\phantomsection\label{ep-003.II}}
\addcontentsline{toc}{subsection}{{\sen{-II-}} {\caviar Mùa hoa phấn của Roboco}}
\stpt{-II-}

\stcap{Mùa hoa phấn của Roboco}
\begin{center}
    \textbf{Kỳ trước ở \hll !}
\end{center}

\pov{Roboco}

Câu chuyện quay ngược trở lại ngày hôm trước, khi tôi nổi hứng đi dạo sâu vào trong rừng mà hoàn toàn không hề hay biết rằng lúc này đang là đỉnh điểm của mùa thụ phấn. Những hạt phấn hoa khổng lồ bay phảng phất khắp không gian, chúng to đến mức bất kỳ ai cũng có thể dễ dàng nhìn thấy bằng mắt thường.

\img{0-4d2}

Và hậu quả là giờ đây, tại văn phòng làm việc thân quen, tôi đang phải ngồi tựa lưng vào ghế với một vẻ mặt vô cùng não nề, thỉnh thoảng lại phát ra một tiếng hắt hơi rõ to. Tôi tự lẩm bẩm đầy tuyệt vọng: ``\textit{Dù mình là một cỗ máy tối tân, thế mà vẫn bị viêm mũi dị ứng trái mùa hành hạ\dots}''

Anh Kuon đứng khoanh tay đối diện tôi, nhướng một bên lông mày lên với vẻ mặt cực kỳ khó hiểu: ``Thú thực là anh vẫn không thể nào hiểu nổi tại sao em lại bị dị ứng được. Với cả\dots\ anh thấy tình trạng của em đang chuyển biến khá nặng rồi đấy, em dùng hết sạch nguyên cả mấy hộp khăn giấy rồi còn gì.''

Tôi mếu máo giơ chiếc khăn giấy đã thấm đẫm một thứ dung dịch màu cam đặc quánh lên, bứt rứt giải thích: ``Đó là vì em đang xì ra dầu máy thay vì nước mũi đó, Kuon-senpai! Giúp em với!''

Cậu Tổng tặc lưỡi lắc đầu: ``Chắc là do em quên không đi bảo trì định kỳ nên hệ thống mới bị lỗi chứ gì. Thôi, hôm nay anh đặc cách cho em về sớm nghỉ ngơi đấy!''

Thế nhưng, tôi chỉ có thể thở dài thườn thượt, tỏ vẻ bất lực cùng cực trước căn bệnh quái gở của mình: ``Nếu em mà tự về nhà được thì em đã về từ đời nào rồi! Dầu của em đang bắt đầu loang lổ khắp sàn văn phòng rồi này!'' Dứt lời, tôi lại vội vàng rút thêm một chiếc khăn giấy khác ra và xì mũi, để rồi chiếc khăn lại tiếp tục ướt đẫm cái thứ chất lỏng màu cam kỳ quặc kia.

Chẳng bao lâu sau, toàn bộ mặt sàn văn phòng đã lênh láng dầu máy trơn trượt. Haato sợ hãi phải nhảy tót lên đứng co rúm trên ghế xoay để né tránh những vệt dầu đang ngày càng loang rộng. Gương mặt em ấy lộ rõ sự lo lắng xen lẫn hoảng hốt: ``Trời đất ơi, dầu của Roboco-senpai vương vãi khắp nơi rồi kìa! Giờ chúng ta phải làm sao đây?''

``\textit{Xin lỗi em nhiều nhé, Haato-chan, nhưng hiện tại chị thực sự cũng hết cách rồi!}''

``Anh cũng đang chịu chết giống hệt em đây này,'' cậu Tổng đứng cạnh Haato thở dài đánh thượt, tay đưa lên xoa nắn thái dương với ánh mắt đan xen giữa sự ngạc nhiên và bất lực tột độ. ``Cũng tại Robo-PON rảnh rỗi đi dạo rừng đúng vào cái mùa phấn hoa nên mới ra nông nỗi này đấy. Mà nực cười ở chỗ, đến tận hôm nay anh mới biết là rô-bốt cũng có thể bị dị ứng đấy.''

``\textit{Giờ không phải là lúc để đứng đó bảo em ngốc đâu nha, Kuon-senpai!}''

\pov{-}

``Mà thế quái nào một con rô-bốt lại có thể mắc bệnh viêm mũi như con người được nhỉ?'' Haato tò mò hỏi, đôi mắt chớp chớp không hiểu.

``Chịu thôi, vụ này thì anh xin bó cả chân lẫn tay,'' Kuon nhún vai đáp lại.

Trong lúc đó, Subaru vẫn đang ngã lăn lộn trên chiếc ghế bành, cười sặc sụa trước tình huống căn bệnh kỳ lạ có một không hai của vị tiền bối. Tiếng cười lanh lảnh không ngừng nghỉ của nàng Vịt hòa quyện cùng những tràng hắt xì văng tung tóe dầu máy của Roboco lại càng khiến cho vấn đề trở nên nghiêm trọng và thảm họa hơn. Cậu Tổng chỉ biết toát mồ hôi lạnh, tay xoa trán đầy bất lực: ``Trong cái hoàn cảnh này mà em vẫn còn ôm bụng cười được thì anh cũng đến lạy em rồi đấy, Subaru-chan\dots''

\img{0-4d}

Với những vết dầu đặc quánh ngày càng loang rộng không có dấu hiệu dừng lại, mọi thứ trong căn phòng này đang dần vượt ra khỏi tầm kiểm soát. Roboco đứng giữa trung tâm văn phòng, bất lực thốt lên một mong ước đầy viển vông: ``Ước gì mình có thể thanh trừng và xóa sạch toàn bộ phấn hoa từ quá khứ, hiện tại, tương lai, và lan ra cả khắp vũ trụ này trước khi chúng kịp sinh sôi!''

Kuon lập tức lùi lại một bước, khoanh tay đứng nhìn cô nàng với vẻ mặt vô cùng cảnh giác: ``Này này, dẹp ngay cái ý tưởng diệt chủng phấn hoa đó đi! Đừng có mà bày ra mấy cái trò nguy hiểm làm ảnh hưởng tới không thời gian nữa!''

Trước khi nàng Người máy kịp có bất kỳ phản ứng nào, tiếng nắm đấm cửa chợt vặn mở. Aqua và Matsuri tung tăng bước vào văn phòng.

``Mọi người ơi, hôm nay em có-- Ôi trời! Sàn trơn quá!'' nàng Hầu gái chưa kịp dứt câu chào thì đã đạp ngay trúng vũng dầu máy trơn tuột mà Roboco vô tình xì ra. Cả cơ thể Aqua mất thăng bằng ngã sóng soài xuống sàn, trán đập mạnh một cú điếng người vào mép bàn làm việc.

Aqua ngồi bệt dưới đất, hai tay ôm khư khư lấy đầu, khóe mắt ngấn lệ vì đau đớn. Nàng Cổ động đứng bên cạnh chứng kiến toàn bộ sự việc, ngay lập tức ôm mặt hét lên để phóng đại tình hình: ``Trời ơi! Aqua-chan đập đầu vào mép bàn rồi! Em ấy bị chấn thương sọ não và mất luôn trí nhớ rồi!''

Kuon nghe vậy liền vỗ trán thở dài, giọng nói tràn ngập sự nghi ngờ: ``Anh nghe thấy sai sai thế nào ấy, Matsuri-chan ạ\dots\ Dù có ngã thì chắc cũng không đến mức khoa trương như thế đâu\dots''

Thế nhưng, sau một hồi ngồi ôm đầu trấn tĩnh, nàng Hầu gái tóc tím-hồng từ từ ngẩng mặt lên. Ánh mắt cô trở nên trống rỗng, giọng nói mơ màng cất lên: ``Tôi\dots\ Tôi là ai? Đây là đâu? Mọi người\dots\ mọi người là gì của tôi vậy?''

Cậu Tổng trợn tròn mắt khi nhận ra cô nàng thực sự không còn nhớ bất cứ thứ gì, hoảng hốt kêu lên: ``Lạy Chúa, ẻm mất trí nhớ thật kìa!''

Nhân cơ hội đó, Roboco lập tức tiến lại gần. Khuôn mặt cô nàng Người máy bỗng chốc bừng sáng lên một sự quyết tâm bất thường. Cô nắm chặt lấy đôi tay của Aqua, như thể một ý tưởng vĩ đại vừa lướt qua trong bộ xử lý trung tâm: ``Giá như mình có sức mạnh của một cô gái ma thuật\dots\ mình chắc chắn sẽ ra tay xóa sổ toàn bộ phấn hoa trên thế gian này để không còn ai phải chịu đựng những nỗi đau khổ như thế này nữa!''

Cậu con trai vội vàng lùi lại thêm một bước, vung tay loạn xạ ngăn cô nàng lại: ``Thôi thôi, anh xin em, dẹp cái ảo tưởng đó đi! Cứ làm như em sắp hắc hóa thành trùm cuối phản diện luôn không bằng ấy!''

Matsuri và Kuon đứng chôn chân nhìn đống dầu đang xâm chiếm văn phòng và cô nàng Hầu gái bị cụng đầu. Trong khi đó, Aqua vẫn ngồi đờ đẫn dưới sàn nhà, liên tục thốt ra những câu nói lơ mơ ngày càng phi logic. Còn Roboco thì đứng vươn mình đầy kiêu hãnh, vẻ mặt hùng hồn như thể đang chuẩn bị đọc một bài diễn văn lịch sử về việc phát động chiến tranh chống lại ``kẻ thù không đội trời chung'' mang tên phấn hoa.

\subsection*{\phantomsection\label{ep-003.III}}
\addcontentsline{toc}{subsection}{{\sen{-III-}} {\caviar \sout{Aqua} Yamada mất trí nhớ}}
\stpt{-III-}

\stcap{\sout{Aqua} Yamada mất trí nhớ}
\begin{center}
    \textbf{Kỳ trước ở \hll !}
\end{center}

\pov{Minato Aqua}

Tôi cầm chiếc chổi quét nhà trên tay, hé nở một nụ cười rạng rỡ và tự tin nhắm nghiền hai mắt lại. Cùng lúc đó, tiếng thuyết minh hào hứng của chính tôi vang lên văng vẳng: ``Trong kỳ trước của \hll!'' Tôi cảm thấy vô cùng tự hào khi có được cơ hội vàng để giới thiệu lại bản thân mình: ``Mình là Yamada, và đây là lần đầu tiên mình đặt chân đến nơi này!''

Tôi bắt đầu cần mẫn lau dọn từng chiếc kệ tủ, rồi cẩn thận lau sạch bóng những chiếc bàn, cố gắng làm việc chăm chỉ nhất có thể. Dẫu rằng tôi đã mất đi toàn bộ ký ức trước đây, nhưng những con người tốt bụng trong văn phòng này đã giúp đỡ tôi rất nhiều!

\begin{figure}[H]
    \centering
    \begin{subfigure}[b]{0.45\textwidth}
        \centering
        \includegraphics[width=8cm]{../../../../Image/Book0/0-4e1.png}
    \end{subfigure}
    \quad
    \begin{subfigure}[b]{0.45\textwidth}
        \centering
        \includegraphics[width=8cm]{../../../../Image/Book0/0-4e2.png}
    \end{subfigure}
    \begin{subfigure}[b]{0.45\textwidth}
        \centering
        \includegraphics[width=8cm]{../../../../Image/Book0/0-4e3.png}
    \end{subfigure}
\end{figure}

Từ đằng xa, tôi nhìn thấy hai người họ đang ngồi thoải mái trên chiếc ghế sô-pha, một nam, một nữ, đang nói chuyện rôm rả với nhau về một vấn đề gì đó. Cô gái tóc tết hai bên bỗng đi về phía tôi, vui vẻ vẫy tay trước mặt tôi và nói một câu gì đó mà tôi nghe không được rõ ràng cho lắm. Dù không hiểu cô ấy đang muốn truyền đạt điều gì, nhưng tôi dám cá chắc chắn rằng cô ấy đang muốn làm quen với một người mới tới như tôi đây mà! Trông gương mặt cô ấy hồi hộp và hào hứng chưa kìa!

\gif{0-4e}{1}{37}

Tôi hăng hái tiến đến chỗ cửa sổ của văn phòng, dùng hết sức mở banh ra để đón lấy những tia nắng ấm áp. Hôm nay đúng là một ngày đẹp trời quá đỗi lý tưởng để làm công việc của một cô hầu gái tận tụy! Có vẻ như tôi đã trót nợ ân tình của họ rất nhiều, nên tôi đang dốc hết toàn bộ sức lực để làm việc chăm chỉ nhằm đền đáp lại công ơn ấy!

Bắt đầu với\dots\ cô nàng tóc ngắn tomboy đang ngồi trên chiếc ghế sô-pha đằng kia! Để tôi chạy tới thử xem cô ấy có cần tôi giúp đỡ việc gì không nào\ldots

Mọi người hãy cùng chống mắt lên xem tôi làm việc hiệu quả và chuyên nghiệp đến mức nào nha!

\pov{-}

Kuon đứng đó đổ mồ hôi hột ròng ròng khi chứng kiến cô nàng ngốc nghếch kia đang chủ động lau dọn sạch sẽ cả văn phòng, một hiện tượng tâm linh kỳ bí mà cô CHƯA từng bao giờ làm kể từ khi bước chân vào làm việc tại công ty này. ``Ít nhất thì dáng vẻ bây giờ em ấy trông cũng giống hệt một cô hầu gái thực thụ rồi, cơ mà\dots''

``\dots Nếu cứ để tình trạng mất trí nhớ này tiếp diễn thì cũng khó khăn thật nhỉ,'' Matsuri nhíu mày nói với vẻ lo ngại. Cô từ từ tiến tới chỗ Aqua, vẫy vẫy tay trước mặt cô nàng và gặng hỏi: ``Này\~{} Akutan ơi! Thật sự không nhớ chút xíu gì về bọn chị luôn sao?''

Nhưng đáp lại cô nàng Cổ động chỉ là một sự im lặng tuyệt đối. Nụ cười công nghiệp vẫn đang nở rạng rỡ trên môi nàng Hầu gái, còn đôi tay thì vẫn đang mải miết chà lau chiếc tủ kính. Lau xong, cô đột ngột quay gót đi thẳng về phía cửa sổ và mở tung nó ra, không hề mảy may bận tâm tới ánh nhìn của những người xung quanh.

``Mắt thì nhắm tịt lại thế kia thì làm sao mà thấy được cái quái gì cơ chứ,'' cậu Tổng thở dài thườn thượt đầy ngao ngán, ``chán thật đấy, hết Subaru-chan, rồi đến Robo-PON, giờ lại lòi ra thêm ca chấn thương sọ não này nữa.'' Cậu nhíu mày quan sát nàng ``Hành tím'' đang lon ton chạy tới hỏi chuyện Subaru xem liệu cô nàng tomboy có cần giúp gì không. Nhưng rồi nàng Vịt chỉ bật cười khanh khách, vẩy tay xua đi và đáp gọn lỏn: ``Bây giờ thì chị chưa cần đâu nghen.''

Nhìn cảnh tượng lạ lùng đến mức vô lý đó, Kuon chỉ có thể bất lực lắc đầu: ``Anh nghĩ là ở đây không cần ai phải trả ơn em đâu, Aqua-chan. Em cứ làm việc chăm chỉ đột xuất thế này khiến anh càng thấy lo sợ hơn đấy.''

Subaru đứng gần đó không thể kiềm chế nổi căn bệnh quái ác của mình, lại bất giác phá lên cười ngặt nghẽo. Cậu quay sang lườm nàng Vịt một cái cháy máy: ``Subaru-chan, làm ơn đừng có cười nữa! Em chỉ đang đổ thêm dầu vào lửa khiến cho mọi chuyện trở nên hỗn loạn hơn thôi!''

Nhưng lời khuyên vừa dứt, nàng tóc ngắn lại càng cười một cách mãnh liệt và điên rồ hơn, đến mức lăn lộn bò trườn luôn dưới sàn nhà. Cậu con trai chép miệng tặc lưỡi: ``Chả biết là em ấy còn đang cười vì cái quỷ quái thánh thần gì nữa đây\dots''

Ngay lúc đó, Ayame - người nãy giờ vẫn đang nằm ườn phè phỡn trên chiếc ghế bành - bất chợt cất lên một giọng nói lười biếng và ngái ngủ: ``Nãy có ai đó nói nhờ vả đúng không ta?'' Cô nàng ngáp dài một cái: ``Tôi nhớ không lầm thì tôi có nhờ bà đi mua cho tôi cái bánh bao thịt nhỉ?'' Cậu Tổng đứng gần đó lập tức chống hai tay ngang hông, phản ứng lại bằng một giọng điệu trách móc: ``Có chân thì tự lết đi mà mua chứ. Em ấy có phải là hầu gái riêng chuyên phục vụ của em đâu--''

Tuy nhiên, trước khi Kuon kịp hoàn thành câu mắng mỏ, Aqua đã hăng hái đập hai tay vào nhau đáp lời: ``Đúng rồi ha! Giờ mình đi mua đồ ăn trưa cho những người đáng mến này nào!'' Nói rồi, cô nàng lập tức xoay người chạy tót ra ngoài cửa như một cơn lốc, nhanh đến mức không ai kịp dang tay ngăn lại.

Cậu hoàn toàn suy sụp, đưa tay lên vò đầu bứt tóc nhìn bóng lưng cô nàng đang chuẩn bị khuất sau cánh cửa. Trong giây phút tuyệt vọng cuối cùng, cậu cố cất lên bằng một giọng chán nản tột độ: ``\dots Phiền em mua luôn cho anh một cốc mì hải sản nhé.''

Aqua gật đầu lia lịa liên hồi, hoàn toàn không nhận ra sự bất mãn và tuyệt vọng trong giọng điệu của Kuon, rồi nhanh chóng vọt đi mất dạng với một nụ cười tươi rói. Cậu gục ngã ngồi phịch xuống ghế ôm mặt mệt mỏi, còn cô nàng tiểu thư Ayame thì vẫn nằm ườn thoải mái trên sô-pha, miệng mỉm cười đắc ý như thể mọi thứ trên đời đang diễn ra đúng theo sự sắp đặt hoàn hảo của mình.

\subsection*{\phantomsection\label{ep-003.IV}}
\addcontentsline{toc}{subsection}{{\sen{-IV-}} {\caviar Giải pháp của Fubuki}}
\stpt{-IV-}

\stcap{Giải pháp của Fubuki}
\begin{center}
    \textbf{Kỳ trước ở \hll !}
\end{center}

\pov{Shirakami Fubuki}

Câu thoại ``Trong kỳ trước của \hll!'' quen thuộc bằng một cách vi diệu nào đó lại do chính tôi lồng tiếng được chèn vào.

Tôi thong thả bước vào văn phòng, nhưng rồi khuôn mặt tôi ngay lập tức cứng đờ, lộ rõ sự ngạc nhiên tột độ khi thấy căn phòng làm việc thường ngày giờ đây lại đang ngập ngụa trong một thứ dầu máy đặc quánh đến tận mắt cá chân. Tôi thấy anh Tổng của chúng ta đang ngồi bó gối thẫn thờ, bất lực đảo mắt nhìn quanh khung cảnh hoang tàn này. Chưa kể, lúc đi dạo vào tôi còn vô tình lướt ngang qua Akutan đang vui vẻ tung tăng chạy ra ngoài. Nhưng cái cách em ấy cười nói vui vẻ lại không hề bình thường chút nào, trông em ấy cứ như đang bị ai đó tẩy não và thao túng tâm lý vậy.

Ngoại trừ anh Kuon vẫn giữ được sự tỉnh táo (dù anh ấy đang trầm cảm), thì hôm nay rốt cuộc mọi người trong công ty bị làm sao thế này? Mà ngẫm lại chẳng phải, mấy ngày gần đây họ đã bắt đầu có những biểu hiện cư xử rất kỳ lạ rồi sao?

Tôi rón rén, từ tốn bước vào sâu hơn, cố gắng không để bản thân bị trượt ngã, khẽ đảo mắt quan sát quanh phòng. Subaru-chan thì đang đứng ở một góc, ôm bụng cười khúc khích đến chảy cả nước mắt mà không rõ nguyên nhân. Ở góc bên kia, Haato, Matsuri, và Ayame đang tụm năm tụm ba bu quanh Aqua, người đang đứng giữa phòng với vẻ mặt lơ ngơ như bò đeo nơ.

``Vậy tên thật của mình là Yamada Hermione sao?'' Aqua chớp chớp mắt, ngây ngô hỏi lại bằng một chất giọng vô cùng hoang mang.

``Không hề có chuyện đó đâu em ơi. Mấy đứa kia lại xúi bậy em cái gì thế?'' Kuon khoanh tay, chép miệng phủ nhận ngay lập tức trước khi ẻm thật sự tin vào cái tên lai Tây quái đản đó.

Phía bên kia sô-pha, chị Roboco đang ngồi ủ rũ trên chiếc ghế bành, một tay cầm chặt cuộn khăn giấy, liên tục hắt xì mạnh bạo như thể cảm giác dị ứng khủng khiếp đang ra sức hành hạ cô ấy. Mỗi lần cô hắt xì như vậy, vài giọt dầu máy nhỏ lại không tự chủ được vương vãi bắn ra sàn, khiến tôi vô thức phải bước lùi lại vài bước để tránh dẫm phải bãi mìn đó. Nhưng điều làm tôi cảm thấy khó hiểu và đáng ngờ nhất là ngoài anh Tổng đang bất lực kia ra, thì còn lại không một ai ở đây cảm thấy lo sợ hay nguy hiểm gì cả.

\begin{figure}[H]
    \centering
    \begin{subfigure}[b]{0.45\textwidth}
        \centering
        \includegraphics[width=8cm]{../../../../Image/Book0/0-4e4.png}
    \end{subfigure}
    \quad
    \begin{subfigure}[b]{0.45\textwidth}
        \centering
        \includegraphics[width=8cm]{../../../../Image/Book0/0-4e5.png}
    \end{subfigure}
    \begin{subfigure}[b]{0.45\textwidth}
        \centering
        \includegraphics[width=8cm]{../../../../Image/Book0/0-4e6.png}
    \end{subfigure}
\end{figure}

Tôi ngước mắt lên, thấy Ayame đang ngồi cười khúc khích ở một góc, còn Haato thì vẫn bình thản ngồi vắt chéo chân trên ghế xoay. Ở góc bếp nhỏ của văn phòng, Aqua vẫn luôn nở một nụ cười tươi rói trên môi, đôi tay thoăn thoắt nhưng có phần lóng ngóng đang rửa một chiếc cốc, hoàn toàn không hề hay biết đến chuyện động trời gì đang xảy ra xung quanh bản thân mình.

Chẳng lẽ chuyện Aqua bị đập đầu mất trí nhớ và chuyện Roboco hắt hơi ra dầu lại thực sự lây lan và ảnh hưởng đến nhận thức của tất cả mọi người đến mức này sao? Hôm nay có lẽ thực sự là một ngày khá bấn loạn và thảm họa đây. Giờ tôi chỉ có thể tự vòng tay ôm lấy mình và cố gắng tự động viên bản thân phải giữ được cái đầu lạnh.

Từ trước đến nay, chúng tôi đã luôn tin tưởng rằng nếu tất cả đồng lòng làm việc cùng nhau, chúng tôi có thể dễ dàng vượt qua bao khó khăn sóng gió. Nhưng nhìn tình cảnh bây giờ thì tôi không dám chắc nữa\dots\ Đến cả một người lý trí như anh Kuon cũng thực sự giơ tay đầu hàng bất lực với cảnh tượng này cơ mà.

Cứ tưởng rằng mọi tia hy vọng đã hoàn toàn vụt tắt, thì bất chợt một ý tưởng lóe sáng trong đầu tôi, khiến đôi tai cáo trên đỉnh đầu tôi vểnh hẳn lên: ``Phải ha! Tại sao mình lại không nghĩ ra cái này sớm hơn nhỉ!'' Đúng lúc này, Matsuri đứng gần đó lên tiếng than vãn: ``Chúng ta thực sự phải làm gì với cái đống dầu bầy hầy này đây? Cứ để yên như thế này thì chả mấy chốc chúng ta sẽ gặp rắc rối to mất!''

Ayame khẽ nghiêng đầu sang một bên, chớp mắt tỏ vẻ không hiểu chuyện gì đang xảy ra.

``Dầu? Dầu gì cơ?'' Ayame ngơ ngác hỏi.

``Thì là đống dầu đang chảy lênh láng thải ra từ mũi của Roboco-senpai đấy!'' Matsuri chỉ tay xuống sàn nhà, lớn giọng giải thích.

``Tức là sự tình nó thế này\dots\ Mà khoan đã, Ojou, em thực sự từ nãy đến giờ không hiểu chuyện quái gì đang xảy ra trong cái văn phòng này sao?'' Kuon trợn tròn mắt, kinh ngạc quay sang nhìn cô tiểu thư quỷ.

Ayame khẽ lấy tay che miệng cười khì khì một tiếng ngộ nghĩnh: ``Hì hì, khum bít nha!''

``\dots Anh xin bó tay toàn tập với em luôn.'' Cậu Tổng xoa xoa vầng trán đang giật giật của mình.

Sau một lúc ổn định lại tinh thần, anh ấy bắt đầu tường thuật ngắn gọn cho cả đám nghe về việc Roboco xui xẻo bị dị ứng phấn hoa trái mùa, và kết quả nhãn tiền là căn phòng thân thương của họ giờ đã biến thành một khu mỏ khai thác dầu thu nhỏ. Khi nghe xong ngọn ngành câu chuyện, mọi người xung quanh mới ồ lên gật gù tỏ vẻ đã hiểu, rồi nhanh chóng tụm lại bắt đầu thảo luận tìm cách xử lý hậu quả.

``Vừa nãy lúc đang nhức đầu thì anh có nghe thấy Fubu-chan reo lên gì đó thì phải nhỉ?'' Kuon đột nhiên cất tiếng hỏi, dường như anh ấy vẫn để ý tới biểu cảm hớn hở của tôi ban nãy. ``Em đang có cao kiến gì sao, cáo nhỏ?''

Tôi vỗ ngực cái bụp, gật gù tự hào, tỏ vẻ vô cùng tự tin với kế hoạch hoàn hảo của mình: ``Đúng á! Mọi người cứ tin ở em, cách này đảm bảo chắc chắn sẽ hiệu quả luôn!''

\ct

Sau một hồi loay hoay hì hục lắp ráp và đóng búa, cuối cùng chúng tôi cũng cùng nhau tạo ra được một tấm bảng lớn có màu trắng tinh khôi. Nếu các bạn đang nghĩ đây chỉ là một tấm bảng viết thông thường, thì các bạn đã nhầm to rồi nhé. Đây chính là tấm bảng siêu thấm hút dầu độc quyền do Fubuki tôi thiết kế!

\img{0-4e7}

Nhưng khổ nỗi, đến lúc chúng tôi làm xong thì lượng dầu trong văn phòng đã dâng cao ngập đến tận đầu gối rồi. Chúng ta phải hành động ngay lập tức thôi. ``Đến giờ dọn dẹp vệ sinh rồi mọi người ơi!'' Tôi hào hứng hô vang.

``Nhưng vấn đề nhức nhối ở đây là, ai sẽ là người đứng ra làm cái chuyện nặng nhọc đó đây?'' nàng Quỷ sừng chống tay hỏi đầy thách thức. Nghe vậy, tôi ngay lập tức nở một nụ cười ranh mãnh, từ từ liếc ánh mắt đầy ẩn ý sang phía Kuon-senpai. ``Tất nhiên ứng cử viên sáng giá nhất chỉ có thể là Kuon-senpai rồi, nhỉ\~{}''

``\dots Ơ khoan đã, cái gì cơ!? Sao lại là anh!?'' anh thảng thốt lùi lại, mặt cắt không còn một giọt máu. Tôi cũng không biết giải thích sao nữa, có lẽ đơn giản vì anh ấy là người đàn ông duy nhất và cũng là người có trách nhiệm nhất ở đây chăng? Tiếng thở dài não nuột của chính anh ấy vang lên hòa quyện cùng những tràng tiếng cười khúc khích trêu chọc của chúng tôi vang dội khắp cả căn phòng.

\subsection*{\phantomsection\label{ep-003.V}}
\addcontentsline{toc}{subsection}{{\sen{-V-}} {\caviar Căn phòng ngập dầu}}
\stpt{-V-}

\stcap{Căn phòng ngập dầu}
\begin{center}
    \textbf{Kỳ trước ở \hll !}
\end{center}

\pov{Natsuiro Matsuri}

Tôi vừa vặn nắm đấm mở cánh cửa văn phòng ra, và ngay lập tức, một dòng thủy triều dầu máy cuồn cuộn tràn vào như thác lũ vỡ đê.

Tôi chỉ biết thở dài thườn thượt khi trơ mắt đứng nhìn dòng dầu đặc quánh không ngừng tuôn trào, rồi quay mặt lại nhìn cả nhóm đang nổi lềnh bềnh: ``Một tấm bảng thấm dầu độc quyền của Fubuki là hoàn toàn không đủ đâu, giờ thì cả văn phòng đã bị nhấn chìm trong biển dầu mất rồi!'' Còn nếu bạn thắc mắc tại sao lượng dầu khổng lồ này lại trào ngược từ bên ngoài hành lang vào, thì đó là vì Roboco-senpai ban nãy đã lon ton chạy ra ngoài để\dots\ hắt hơi cho thoải mái, dẫn đến việc một lượng lớn dầu máy bị tích tụ ngoài đó. Và rồi khi cánh cửa vừa hé mở, chúng lập tức tràn vào như thể đập thủy điện xả lũ, cuốn trôi theo cả chính cô nàng ``tác giả'' của đống dầu thảm họa này.

\begin{figure}[H]
    \centering
    \begin{subfigure}[b]{0.45\textwidth}
        \centering
        \includegraphics[width=8cm]{../../../../Image/Book0/0-4f1.png}
    \end{subfigure}
    \quad
    \begin{subfigure}[b]{0.45\textwidth}
        \centering
        \includegraphics[width=8cm]{../../../../Image/Book0/0-4f2.png}
    \end{subfigure}

    \vspace{0.5cm}

    \begin{subfigure}[b]{0.45\textwidth}
        \centering
        \includegraphics[width=8cm]{../../../../Image/Book0/0-4f3.png}
    \end{subfigure}
    \quad
    \begin{subfigure}[b]{0.45\textwidth}
        \centering
        \includegraphics[width=8cm]{../../../../Image/Book0/0-4f4.png}
    \end{subfigure}
    \caption{Căn phòng ngập dầu với biểu cảm và hành động ``ai ở đâu thì ở yên đó'' của các nhân vật.}
\end{figure}

Trong phòng, tất cả mọi người đang dần nổi lơ lửng trên một mặt hồ dầu trơn trượt. Chị Roboco ngồi gác chân trên ghế bành trôi dạt, Ayame thì thản nhiên khoanh tay nhắm mắt tịnh tâm, Haato vẫn kiên quyết bám trụ trên chiếc ghế xoay đang xoay vòng vòng, Aqua vẫn tiếp tục duy trì nụ cười mơ màng ngốc nghếch, còn Subaru\dots\ thì đang chật vật cố gắng giữ thăng bằng trong biển dầu mà vẫn không thể ngừng được những tràng cười sặc sụa. Dù tình thế hiện tại khiến mọi người không thể há miệng trò chuyện trọn vẹn, cả nhóm vẫn xuất sắc giao tiếp với nhau bằng những bong bóng dầu nổi lên sùng sục và ánh mắt đưa tình đầy ý nhị, tạo ra một bầu không khí ``quá lố'' đến mức nực cười.

Đến cả vị Tổng quản lý đáng kính Kuon-senpai cũng không thể nào tránh khỏi số phận bi đát ấy. Anh thậm chí còn chán nản đến mức không buồn vùng vẫy bơi lội nữa, chỉ buông thả cơ thể trôi nổi tự do rồi lắc đầu đầy ngao ngán: ``Đến cái nước này rồi mà mọi người vẫn còn có thể cười nói vui vẻ được thì con cũng xin ạ\dots''

Tôi đưa mắt đảo quanh một vòng, cố gắng hít một hơi thật sâu để đưa ra một quyết định mang tính lịch sử: ``Và cuối cùng, chúng tôi sẽ quyết định di dời văn phòng đi chỗ khác!'' Tôi lờ mờ nhìn thấy anh Tổng đang muốn giơ tay lên phản đối quyết định bốc đồng này, nhưng sau cùng anh ấy cũng chỉ biết bất lực lắc đầu, ánh mắt như muốn gào thét: ``\textit{Anh cứ tưởng mấy đứa sẽ có cách nào giải quyết thiết thực và bình thường hơn cơ đấy\dots}''

\subsection*{\phantomsection\label{ep-003.VI}}
\addcontentsline{toc}{subsection}{{\sen{-VI-}} {\caviar Bếp lửa}}
\stpt{-VI-}

\stcap{Bếp lửa}
\begin{center}
    \textbf{Kỳ trước ở \hll !}
\end{center}

\pov{Nakiri Ayame}

Tôi nằm dài trước màn hình máy tính, hai tay ôm lấy chiếc bụng đang sôi lên ùng ục vì đói lả. Có lẽ là do sáng nay tôi đã nướng khét lẹt trên giường đến tận trưa trật trưa trờ, nên tôi vẫn chưa có bất kỳ thứ gì lót dạ cả, ngoại trừ mấy món đồ ăn vặt linh tinh mà Aqua mang về từ đời thuở nào.

Tôi cố gắng nghển cổ ngó sang xung quanh và chợt nhận ra rằng cô nàng Hầu gái ngốc nghếch kia nãy giờ vẫn chưa chịu quay lại từ sau nhiệm vụ đi mua đồ ăn trưa. Bằng một nghị lực phi thường nào đó, con bé đã xuất sắc bơi lội thoát được khỏi cái văn phòng ngập ngụa dầu máy này để ra ngoài, nhưng tôi chờ mòn mỏi mãi vẫn chưa thấy nó mang bánh bao thịt về, mà bụng tôi thì đã đánh trống kêu ca cồn cào lắm rồi.

``\textit{Nếu con bé mà không về nhanh, chắc mình sẽ phải tự thân vận động lăn xả vào trong bếp thôi nhỉ?}''

Thở dài thườn thượt, tôi quyết định tự bơi vào bếp để xem có gì nấu ăn lót dạ hôm nay hay không. Thế nhưng, ngay khi tôi vừa mới nổi lềnh bềnh đến sát chiếc bếp ga, Matsuri và Fubuki bất ngờ từ đâu ``bơi'' xộc tới. Cả hai người bọn họ giơ hai tay quơ loạn xạ trong không trung và dưới mặt dầu. Bọn họ đang cố gắng nhô đầu lên phát ra âm thanh gì đó giữa cái tình trạng ngập ngụa dầu máy lên tận cổ này, nhưng tiếc thay, thứ duy nhất lọt vào tai tôi chỉ là những tiếng bọt khí nổ ``ọc ọc'' đầy khó hiểu.

\img{0-4g}

Ở cách đó không xa, anh Tổng con người của chúng tôi lén lút liếc nhìn qua cảnh tượng hỗn loạn đó, rồi vội vã tìm cách trượt qua khe cửa và chuồn êm khỏi căn phòng. Có chuyện gì xảy ra với anh ấy vậy nhỉ? Sao tự dưng lại bỏ chạy nhát cáy thế?

``\textit{Thôi kệ đi. Chắc anh ấy mắc đi vệ sinh hay là có việc gì đó đột xuất thôi.}'' Tôi nhún vai quay đầu lại nhìn về phía chiếc bếp, rồi lại hướng mắt về phía những bong bóng dầu đang sủi bọt ùng ục, nơi hai người bạn chí cốt của tôi đang tuyệt vọng cố gắng truyền đạt với tôi một thông điệp sinh tử nào đó. Họ đang cố gắng cảnh báo tôi điều gì đó sao?

Nhưng thực sự thì với thứ ngôn ngữ bong bóng này, tôi cũng xin chịu thua thôi!

Bất chấp tình hình ngập lụt hỗn loạn, tôi vẫn ung dung tiếp tục công việc nội trợ của mình. Tôi cẩn thận nhấc chiếc chảo đặt lên trên chiếc bếp ga. \textit{Chắc tôi cũng nên xào ít rau hay làm món gì đó ăn tạm cho qua bữa thôi nhỉ.}

Giờ thì, chỉ cần vặn núm bật bếp ga lên và\dots

\pov{-}

Ngay khi Nakiri Ayame vừa xoay núm bật chiếc bếp ga, một tia lửa mồi rực rỡ bắt đầu lóe sáng lên. Và tất nhiên, ở một nơi ngập tràn khí gas và dầu máy đặc quánh xung quanh, tia lửa nhỏ bé ấy chính là lời tiên tri báo hiệu cho một sự cố nổ tung trời lở đất sắp diễn ra. Ở góc dưới bên trái màn hình, một mũi tên vàng khè mang dòng chữ \textit{To be continued} lừng lững hiện ra. Đồng thời, giai điệu bass điệu nghệ của bài hát \href{https://www.youtube.com/watch?v=cPCLFtxpadE&t=40s}{Roundabout của ban nhạc Yes} bỗng văng vẳng cất lên, khép lại phân cảnh này một cách vô cùng nghệ thuật và đầy kịch tính.

\img{0-4h}

Chắc các bạn cũng đoán được cái kết sau đó rồi. Đây sẽ là một cú nổ lớn đấy.
%==========================================TRUYỆN PHỤ=========================================
\omk

\subsection*{\phantomsection\label{ep-003.VII}}
\addcontentsline{toc}{subsection}{{\sen{-VII-}} {\caviar Cuộc hỗn loạn bùng cháy} (Good Ending)}
\stpt{-VII-}

\stcap{Cuộc hỗn loạn bùng cháy}
\begin{center}
    \textbf{Kỳ trước ở \hll !}
\end{center}

Ayame vừa vặn núm bật chiếc bếp ga lên, không ai ngờ tới rằng toàn bộ lượng dầu máy đặc quánh trong phòng lập tức bắt lửa chỉ trong chớp mắt. Thế nhưng, thay vì tạo ra một vụ nổ kinh hoàng banh xác giống như lẽ thường tình, ngọn lửa lại bùng cháy lên một cách chậm rãi nhưng đầy mãnh liệt, lan rộng ra khắp bề mặt dầu tạo thành một cảnh tượng rực lửa y như cắt ra từ một bộ phim thảm họa Hollywood. Fubuki và Matsuri hốt hoảng quẫy đạp, cố gắng vẫy vùng dưới mặt biển dầu đang rực cháy để ra tín hiệu cầu cứu, nhưng vẫn không làm cách nào hé miệng hét lên được những tiếng kêu cứu tử tế.

Chả mấy chốc, toàn bộ không gian văn phòng đã bị bao trùm trong biển lửa đỏ rực, ngọn lửa hung hãn bốc lên cao chặn đứng hẳn lối ra vào khiến không ai có thể bơi thoát ra ngoài được. Tất cả bọn họ cũng không thể hét lên cầu cứu được vì sợ sặc dầu. Cứ tưởng rằng kết cục bi thảm nhất là tất cả mọi người sẽ bị chết cháy trong thảm họa này\dots

``\dots Ủa, dòng dầu đang từ từ rút xuống kìa mọi người!'' Ayame, người nãy giờ vẫn đang ngụp lặn trên bề mặt biển dầu bốc cháy, thảng thốt lên tiếng trong sự ngạc nhiên tột độ. Cô nàng không biết là do họ đang gặp may, hay là vì một phép màu tâm linh nào đó đang xảy ra. Mức dầu cứ thế chảy chậm dần, chậm dần, rút đi đi qua ngưỡng cửa sổ vỡ, cho đến khi chỉ còn ngập đến ngang đầu gối.

Roboco bỗng dưng ngoi lên xuất hiện ngay bên cạnh Ayame, hai đường vệt dầu vẫn đang chảy tong tỏng ra từ mũi, nhưng rõ ràng là không còn phun trào mạnh mẽ như trước nữa. Cô nàng Người máy cất tiếng tuyên bố với một sự quyết tâm cao độ: ``Mọi người tránh ra! Để chị dập tắt đám cháy này!'' như thể cô đang muốn tự tay chuộc lại lỗi lầm thảm họa do chính mình gây ra.

Nói là làm, nàng Rô-bốt ngay lập tức kích hoạt chế độ cứu hỏa khẩn cấp. Từ hai lòng bàn tay cô, hai luồng nước áp lực cao bắt đầu phun ra xối xả. Nhưng trời tính không bằng Roboco tính, những tia nước ấy khi bắn vào biển dầu đang cháy lại vô tình khiến dầu văng tung tóe khắp nơi, trực tiếp làm đám cháy lan rộng và bốc cao hơn nữa. Nàng Quỷ sừng chỉ biết đứng nhìn khung cảnh rực lửa ngày càng tồi tệ đó mà không sao làm gì để ngăn cản được. ``\textit{Quả thật Kuon-senpai nói cấm có sai\dots\ Roboco-senpai đúng là chúa PON mà.}''

Bất chợt, Aqua từ đâu ngụp lặn trồi lên cùng với hai người họ, trên người còn đang tự trang bị một chiếc phao bơi cứu sinh màu hồng phấn. Những thành viên còn lại cũng đã lần lượt ngoi lên khỏi mặt dầu, ai nấy đều tỏ ra hoảng loạn và nhếch nhác. Cô nàng Hầu gái điềm tĩnh vỗ ngực trấn an tất cả mọi người: ``M-M-Mọi người à! Đừng lo lắng, cứ để tôi lo liệu\dots\ Ê-tồ\dots\ Tôi\dots''

Tất cả mọi người đều đồng loạt nín thở, hồi hộp chờ đợi câu trả lời và kế hoạch thần thánh của cô. Nhưng rồi cô nàng chỉ lơ ngơ nghiêng đầu, thốt lên một câu nói vô tri đầy trống rỗng: ``Nãy tôi tính làm gì ấy nhỉ?''

Đúng lúc đó, Kuon hộc tốc chạy quay lại văn phòng với một khuôn mặt kinh hoàng tột độ khi chứng kiến một bãi chiến trường rực lửa đỏ rực đang thiêu rụi công ty! ``Chuyện quái quỷ gì đang diễn ra ở đây thế này!?'' Cậu thậm chí còn cảm nhận được hơi nóng hầm hập của ngọn lửa và mùi khét lẹt của dầu máy bốc ra từ bên ngoài hành lang.

Cậu Tổng đảo mắt nhìn về phía cô nàng Rô-bốt vẫn đang ra sức xịt nước vào đống lửa, nhưng càng dập, lửa lại càng cháy bừng bừng bạo liệt hơn. ``À-rế, tại sao lửa lại không chịu tắt vậy ta?'' Roboco ngây thơ thắc mắc.

Câu nói vô tư của Roboco khiến cho tất cả mọi người có mặt trong phòng đều phải ngán ngẩm đồng thanh kêu trời: ``\textbf{Tôi chưa từng được học qua khóa đào tạo nào bảo phải dập lửa bằng cách xịt nước như vậy đâu nha!}'' Haato gào lên.

Fubuki, Matsuri, Ayame và Aqua thì đồng thanh thét lớn: ``\textbf{Roboco-senpai đúng là đồ ngốc nghếch mà!}''

``Lạy chúa tôi ơi, ngu đến thế là cùng! \textbf{Ai lại đi dùng nước để dập lửa cháy dầu cơ chứ!?}'' Kuon ôm đầu tuyệt vọng gào thét.

``Ể? Thế không phải là lửa nào cũng có thể bị nước dập tắt sao?'' Roboco chớp chớp mắt vô tội.

Có lẽ vì cậu chàng Tổng quản lý đã quá kiệt sức và mệt mỏi với những chuyện hoang đường xảy ra trong ngày hôm nay, nên cậu cũng chẳng thèm phí lời để giải thích kiến thức vật lý cơ bản cho cô nàng nữa. Thay vào đó, cậu lập tức tiến lên chỉ huy cả nhóm, thể hiện phong thái của một người gánh vác cả công ty: ``Roboco! Bây giờ thì tắt ngay cái vòi nước đó đi!''

Cô nàng lật đật làm theo y như lời cậu bảo. Kế tiếp, cậu quay sang dặn dò: ``Matsuri-chan! Em chạy ra ngoài lấy ngay cái bình chữa cháy bột cho anh! Nó được treo ở trên tường ngay gần cửa thoát hiểm đấy!'' Nàng Cổ động nghe vậy lập tức bừng tỉnh, đáp lớn: ``Rõ! Vâng ạ!'' rồi lội bì bõm qua đống dầu chạy vọt ra ngoài.

``Fubuki-chan, đập vỡ hết cửa sổ đi để dầu tràn thoát hết ra ngoài!'' Không chần chừ do dự thêm một giây nào, nàng Cáo trắng nhanh nhẹn vận dụng kỹ năng ninja đập vỡ toang các khung kính, trong khi những người còn lại thì cố gắng trèo lên những chỗ cao hơn như kệ tủ để không bị bén lửa. Đứng giữa mớ hỗn độn đó, nàng Người máy vẫn còn ngần ngại thắc mắc: ``Em cứ ngỡ văn phòng mình xịn xò lắm phải có lắp hệ thống chữa cháy tự động trên trần chứ?''

Kuon ngán ngẩm liếc cô nàng, đáp lời một cách đầy chua chát: ``Em nghĩ là với cái công ty nghèo kiết xác đến mức còn chưa được dọn dẹp tử tế thế này thì tiền đâu mà họ lắp hệ thống chữa cháy tự động cho em?'' Roboco chỉ biết há hốc mồm ``Ồ\dots'' một tiếng hiểu chuyện, rồi gật gù đứng nhìn hai người kia hợp sức dập lửa.

Bất chợt, trong khoảnh khắc nước sôi lửa bỏng, nàng Hầu gái mất trí Aqua đột nhiên nảy ra một ý tưởng vô cùng táo bạo. Cô hét lớn đầy dõng dạc: ``Mọi người mau lấy hết mì ly cốc ra đây! Nấu nước sôi đổ lên lửa đi! Lấy độc trị độc!'' Lập tức, tất cả động tác của mọi người đều chững lại, quay sang nhìn cô bằng ánh mắt kỳ thị và khó hiểu nhất có thể, nhưng rồi nàng Quỷ sừng Ayame cũng hùa theo, tặc lưỡi đáp: ``Thôi thì\dots\ đằng nào cũng thế, cứ thử đại đi xem sao!''

\ct

Sự hỗn loạn và hoảng loạn của cả văn phòng kéo dài mãi cho đến khi đám cháy cuối cùng cũng được kiểm soát nhờ chiếc bình chữa cháy của Matsuri. Cả nhóm mệt nhoài ngồi bệt xuống sàn nhà vẫn còn đang nhầy nhụa dầu máy, đồng loạt thở phào nhẹ nhõm nhưng ai nấy cũng đều liếc nhìn Aqua với ánh mắt đầy phán xét và nghi ngờ. Dường như cô ấy vẫn chưa thể khôi phục lại ký ức sau cú va đập lúc nãy. Cô nàng chỉ biết gãi đầu cười trừ, ngây ngô hỏi một câu: ``Vậy là\dots\ chúng ta đã giành chiến thắng trước bọn phấn hoa rồi chứ?''

Tất cả mọi người trong nhóm đều buông một tiếng thở dài thườn thượt đầy chán nản. Còn Kuon thì chỉ biết gục đầu lên mặt bàn làm việc ám khói đen, rên rỉ than thở: ``Đúng là anh xin hàng chịu thua mấy đứa rồi đấy\dots\ Cơ mà nhắc mới nhớ, nãy giờ mọi người có thấy Subaru-chan ở đâu không?'' Các cô gái nghe đến tên nàng Vịt liền giật mình bật dậy, cuống cuồng nhìn ngó dáo dác quanh phòng tìm kiếm.

``Đúng á! Nãy giờ em vẫn nhớ cậu ta còn ở quanh đây cười mà!'' Matsuri hoang mang nói. Fubuki vuốt vuốt mái tóc trắng đang lấm lem khói, gãi đầu lẩm bẩm tự vấn: ``Không biết em ấy bơi đi đâu mất rồi nhỉ?''

Và hóa ra, nàng Vịt tomboy của chúng ta đã cùng với số dầu rò rỉ kia trôi tuột ra ngoài mặt đường qua khung cửa sổ. Suốt dọc chuyến hành trình trôi nổi lênh đênh bị dòng dầu đẩy đi khắp mọi ngóc ngách của thành phố đó, cô nàng vẫn không thể ngừng được căn bệnh ôm bụng cười sặc sụa của mình.

Báo hại cả công ty nguyên ngày hôm đó phải chia nhau bơi đi tìm cô ấy mệt bở hơi tai!

\subsection*{\phantomsection\label{ep-003.VIII}}
\addcontentsline{toc}{subsection}{{\sen{-VIII-}} {\caviar Sụp đổ} (Bad Ending)}
\stpt{-VIII-}

\stcap{Sụp đổ}
\begin{center}
    \textbf{Kỳ trước ở \hll !}
\end{center}

\pov{Hikawa Kuon}

Bằng một cách thức trốn chạy thần kỳ nào đó, cuối cùng tôi cũng may mắn thoát khỏi cái bể chứa dầu tử thần ngột ngạt kia. Ngay sau khi tôi vừa bước chân chạy thoát ra khỏi cánh cửa văn phòng thì\dots

{\Large \textbf{BÙM!}}

Một tiếng nổ đinh tai nhức óc vang lên chấn động cả khu phố. Cả văn phòng đã phát nổ dữ dội và bị ngọn lửa cuồng nộ thiêu rụi hoàn toàn. Trực giác của tôi mách bảo đâu có sai, lửa gặp dầu máy nguyên chất mà không tạo ra một vụ nổ hoành tráng mới là chuyện lạ đó! Tôi chỉ biết đứng đó, tuyệt vọng nhìn chằm chằm vào những đám khói đen kịt đặc quánh đang ùn ùn bốc lên từ phần tàn tích đổ nát còn sót lại của văn phòng làm việc xấu số kia.

Ngước mắt nhìn lên bầu trời xanh vợi, tôi bàng hoàng chứng kiến cảnh các cô đồng nghiệp thân thương của mình đang rơi rụng tự do xuống đất y như những viên thiên thạch nhỏ. Trên người họ vừa bị cháy xém đen thui lại vừa tỏa sáng lấp lánh\dots\ theo một cách vô cùng phản khoa học và khó hiểu. Roboco rơi tự do trúng ngay phóc xuống đầu Aqua - người vẫn đang gồng mình ôm chặt túi đồ ăn vặt lúc nãy mua, khiến cả hai cô nàng nổ đom đóm mắt và bất tỉnh nhân sự tại chỗ.

Haato, Matsuri, Fubuki, và Ayame lồm cồm bò dậy từ đống đổ nát, toàn thân lấm lem bụi than, dầu máy và khói bụi đen xì. Tuy nhiên, ngoại trừ việc trông hơi tàn tạ ra thì không có ai trong số họ có vẻ gì là đang cảm thấy đau đớn cả. Bọn họ vội vã chạy tới xem xét tình trạng của hai cô nàng đen đủi kia. Ở một góc khác, chỉ có mỗi Subaru là vẫn đang nằm lăn lộn ra đất cười nắc nẻ sằng sặc, như thể thảm họa vừa rồi chính là một show hài kịch tuyệt đỉnh nhất mà cô từng được xem.

``Ha ha ha! Phụt! Này Kuon-senpai, anh nhìn xem này, cái bộ dạng bọn này buồn cười chết mất! Ha ha ha!'' Subaru ôm quằn quại dưới đất, chỉ tay vào mấy người kia trêu chọc.

Tôi chỉ có thể chép miệng, buông một tiếng thở dài đứt ruột: ``Ừ thì\dots\ chỉ là hậu quả của một pha nghịch dại đi vào lòng đất thôi mà.''

Tôi lôi chiếc điện thoại di động trong túi quần ra, bấm số gọi cho Hijikara Khan (\ruby{土}{ひじ}\ruby{方}{かた}\ruby{判}{はん}, Đỗ Phương Phán) - người bạn từng học hồi cấp ba với tôi, giờ đang làm bác sĩ trong một bệnh viện tư ở Kawachi.

``\vie{A lô, Khan hả ông? Đặt khẩn cấp giúp tôi hai suất giường bệnh đặc biệt ở bệnh viện với. Một ca chấn thương sọ não mất trí nhớ, một ca tâm thần. Gửi thẳng hóa đơn về công ty tôi nhá, nhờ ông giúp ca này đấy.}''

Cúp máy, tôi khẽ đảo mắt nhìn lại cả nhóm đang chật vật cố gắng kéo lê xác Roboco và Aqua ra khỏi đống xà bần ngổn ngang bốc khói, rồi tôi lại phải lôi điện thoại ra gọi thêm một cuộc điện thoại nữa cho bố mình. ``\vie{Bố ơi, bố có thể chạy xe ra giúp con đưa Roboco-san về lại xưởng cơ khí nhà mình để sửa chữa lại đi ạ. Cô ấy vừa phát nổ và hiện đang bị rò rỉ dầu trầm trọng lắm rồi.}''

Ngắt máy điện thoại, tôi khoanh tay đứng lặng yên, bất lực đưa mắt nhìn cảnh tượng hoang tàn đổ nát và hỗn loạn trước mặt. Trong đầu tôi giờ đây chỉ văng vẳng lặp đi lặp lại một câu hỏi vô tận: ``\textit{Công ty này\dots\ rồi tương lai sẽ đi về đâu đây\dots?}''

Tôi lại thở hắt ra một tiếng thở dài não nề và đầy mệt mỏi. Ánh mắt tôi hướng về một chân trời xa xăm, nghĩ ngợi về cái tương lai xám xịt và không mấy sáng sủa khi quyết định gắn bó lâu dài ở công ty này. Có lẽ, quyết định đồng ý vào làm việc quản lý cho lũ khùng này không phải là một ý kiến hay ho gì.

Nhưng biết làm sao được chứ, vì miếng cơm manh áo, và cũng vì để chứng minh cuộc sống tự lập với gia đình, lại còn làm đúng ngành và lương cao, tôi thực sự cũng chẳng còn con đường nào khác ngoài việc tiếp tục nghiến răng cam chịu và làm việc tại đây.

Hầy da. Đời tôi là vậy đó.

\end{document}
